\section{Discussion}
\label{sec:discussion}

\subsection{AreaSection as a Geographic Fairness Mechanism}

The central empirical finding of this paper is that AreaSection changes the
\emph{structure} of redistricting outcomes without changing the \emph{outcome}
for partisan seat counts.
In both Wisconsin and North Carolina, AreaSection and GeoSection produce
identical seat allocations and identical proportionality gaps.

This finding has an important implication: the proportionality gap in these
states is not a product of urban peeling.
It is a product of the underlying population distribution.
A 50.3\,\% Democratic state (Wisconsin) with Republicans distributed across
a larger geographic footprint will consistently produce a 3D/5R congressional
delegation under any geographically neutral algorithm, regardless of whether
the bisection strategy is urban-peeling (GeoSection) or area-balanced
(AreaSection).

What AreaSection \emph{does} change is the competitiveness profile:
knife-edge districts disappear.
GeoSection's suburban ring districts at 49.4\,\% Democratic vote share are a
systematic artifact of the peeling strategy.
AreaSection's cleaner partisan sorting --- safe Democratic districts and
clearly Republican districts --- is arguably more honest about the underlying
partisan geography.

\paragraph{Legal theory.}
The area balance constraint has a natural legal justification:
census tracts are administrative units whose areas reflect the physical
geography of the state, not the partisan composition of their residents.
A redistricting algorithm that enforces geographic balance respects the
physical structure of the state in a way that population-only algorithms
do not.

Under \emph{Rucho v. Common Cause} (2019), federal courts cannot remedy
partisan gerrymandering under the Constitution, but state courts can and
have done so \citep{lwv2018pa,harper2022nc}.
State courts reviewing redistricting algorithms may find the geographic
balance criterion independently defensible:
``we drew district boundaries to minimise total boundary length while
giving each half of the state comparable geographic footprint''
is a claim with a clear, politically neutral rationale.

\subsection{Why Seats Don't Change}
\label{sec:why-seats-stable}

The stability of seat counts across GeoSection and AreaSection reflects
a deeper mathematical structure.
In states with strong geographic partisan sorting (urban Democrats,
rural Republicans), the \emph{number} of Democratic seats is determined
by how many contiguous geographic regions have Democratic majorities,
not by where the bisection boundaries fall.
AreaSection changes where the boundaries fall, but not which geographic
regions are sufficiently Democratic to elect a Democrat.

This suggests that the partisan outcome is, in a precise sense, a property
of the population distribution, not of the bisection algorithm.
GeoSection (B.7) establishes this claim rigorously for the seed-sweep
convergence case.
AreaSection provides corroborating evidence from a different direction:
changing the geometric constraint does not change the seat count.

\subsection{Limitations}
\label{sec:limitations}

\paragraph{Large-state balance failures.}
Six states --- FL, IL, MI, NY, PA, TX --- fail the 1.5\,\% population balance
check.
The root cause is recursive bisection rounding: when the first-level split
produces a subgraph with an unusual population distribution (e.g., 35\,\%
instead of 37.5\,\% of the state population), the recursive subgraph cannot
balance to the national ideal per-district within the discrete-tract
constraint.

This is not a fundamental limitation of AreaSection but a rounding artefact
that can be addressed by:
(a) increasing the recursive bisection seed count to find better-balanced cuts,
(b) relaxing the per-subgraph tolerance, or
(c) post-processing with a population-balancing swap algorithm.
The 44-state results demonstrate the algorithm's correctness on states where
this artefact does not arise.

\paragraph{Area vs population correlation.}
The area constraint is an indirect measure of geographic diversity.
In states where the area-to-population ratio is approximately uniform
(agricultural Midwest), AreaSection behaves identically to GeoSection.
The constraint only has meaningful effect in states with strong
population concentration, which are precisely the states where urban peeling
is most problematic.

\paragraph{The 10\,\% swing tolerance.}
The $\pm 10\,\%$ area swing (ubvec[1] = 1.10) is a design choice.
Tighter tolerances (e.g., $\pm 5\,\%$) would force more equal geographic splits
but would increase the number of Lorenz-infeasible states.
Looser tolerances (e.g., $\pm 20\,\%$) would permit more urban peeling.
The 10\,\% value is chosen to match the natural geographic variability of
census-tract configurations while still preventing extreme peeling.
Future work should test sensitivity to this parameter.

\paragraph{Area measure.}
We use TIGER ALAND (land area only, excluding water).
ALAND + AWATER would change results for coastal and Great Lakes states
(NY, MI, WI) by attributing large water areas to coastal tracts.
ALAND is the more legally defensible choice: water is not habitable territory
and should not count toward a district's geographic footprint.

\subsection{Comparison to Polsby-Popper and Other Compactness Metrics}

Standard compactness metrics (Polsby-Popper, Reock, convex-hull ratio) evaluate
the shape of individual districts after the fact.
AreaSection is an \emph{algorithmic} constraint applied during bisection:
it shapes the structure of the bisection tree before any individual district
boundaries are drawn.

These are complementary, not competing, approaches.
AreaSection produces bisection trees that respect geographic balance;
Polsby-Popper audits verify that individual districts within those trees
are not pathologically non-compact.
A complete redistricting analysis would apply both.
